\section{分岐自由因子と普通二重点による正規化}\label{sec:four-branch-v24}

前節の四枚被覆について、全空間のramification divisorと基底のbranch divisorを
区別して記述する。九つの自明な変数を除き、
$R=\mathbb C[[u,v,w,z]]$、$\mathscr D\simeq\mathbb C[[X,Y,w,z]]$ とする。
この節の行列式は多項式恒等式であり、有限jetへの打切りを含まない。

\subsection{trace判別式と対数的ベクトル場}

\begin{theorem}\label{thm:four-free-discriminant-v24}
次の行列と多項式を定める。
\begin{equation}\label{eq:four-discriminant-matrix-v24}
C=\begin{pmatrix}
 uw/4&-v/2&2u&(-3uz+vw)/4\\
 3uz&4u-wz&2v&(-uw-vz+w^2z)/2\\
 (8v-w^2)/4&-3z&w&(16u+5wz)/4\\
 -4u&3w/2&z&(4v-3w^2)/4
\end{pmatrix},\qquad \delta=\det C.
\end{equation}
四枚被覆のbranch divisorは被約既約な超曲面 $\Delta=V(\delta)$ である。
$\delta$ は基底 $(1,X,Y,X^2)$ に関するtrace判別式であり、
重み $(2,2,1,1)$ について次数八の準斉次多項式である。
$-C$ の四つの列を係数とするベクトル場は
$\Der_R(-\log\Delta)$ の基底を与える。特に $\Delta$ は自由因子である。
さらに、$q=4X^2+2Y^2+wX-zY$ と置くと
\begin{align}
 \Nm_{\mathscr D/R}(q)&=-8\delta,\label{eq:four-norm-v24}\\
 \operatorname{disc}_X(2X^4+wX^3+vX^2+zuX-u^2)&=16u^2\delta.
 \label{eq:four-quartic-index-v24}
\end{align}
従って四次多項式の判別式には、四枚被覆自体の分岐に由来しない因子 $u^2$ がある。
\end{theorem}

\begin{proof}
$\mathscr D$ の有限自由基底 $(1,X,Y,X^2)$ における乗法行列は
\[
M_X=\begin{pmatrix}
0&0&u&-zu/2\\1&0&0&-v/2\\0&0&0&u/2\\0&1&0&-w/2
\end{pmatrix},\qquad
M_Y=\begin{pmatrix}
0&u&v&0\\0&0&w&u\\1&0&z&0\\0&0&2&0
\end{pmatrix}.
\]
直接乗算により $M_XM_Y=uI$、
$M_Y^2-2M_X^2=vI+wM_X+zM_Y$ が成り立つ。
$M_1=I,M_2=M_X,M_3=M_Y,M_4=M_X^2$ とすると
\[
 \det\bigl(\operatorname{Tr}(M_iM_j)\bigr)_{i,j=1}^4=\det C.
\]
同じ行列の行列式から式\eqref{eq:four-norm-v24}が得られ、
四次多項式のresultantから式\eqref{eq:four-quartic-index-v24}が得られる。
これらは有理数係数の多項式の直接計算である。
trace判別式による非\'etale locusの記述は
\cite[Tag 0BVH]{Stacks}の標準的な判定とも一致する。

行列 $C$ は相対接コホモロジーの最小表示から得られる。
\[
 T^1_{\mathscr D/R}=\coker\left[
 \mathscr D^2\xrightarrow{
 J=\left(\begin{smallmatrix}Y&X\\-4X-w&2Y-z\end{smallmatrix}\right)}
 \mathscr D^2\right]
\]
を上の基底で八次正方行列に展開する。
行 $(2,3,4,8)$、列 $(1,2,3,5)$ の小行列式は $-8$ である。
この単位ブロックを消去したSchur補行列が $C$ であり、
同じ消去の下で $u,v,w,z$ のKodaira--Spencer写像は $-I_4$ になる。

ここで被約性と既約性を確認する。
$V(q)$ は $\alpha\beta+\gamma\eta=0$ という三次元普通二重点なので整かつ正規である。
有限射によるその像は既約であり、normの等式からその台は $V(\delta)$ に等しい。
多項式族の基底点 $(u,v,w,z)=(1,-7,0,6)$ では
\[
 2X^4-7X^2+6X-1=(X-1)^2(2X^2+4X-1).
\]
唯一の重根は $X=1$ であり、対応する点は $(X,Y)=(1,1)$ である。
ここで $\ker J$ は $(1,-1)$ で生成され、
$dq(1,-1)=10$ だから単純なfoldである。
従って像の一般点における判別式の重複度は一である。
形式完備化の後も $V(q)$ は整であり、有限射の像の台は既約である。
また準斉次多項式 $\delta$ の被約性は完備化で失われない。
実際、形式級数の非単元の二乗が $\delta$ を割れば、最低重み項の比較により
多項式 $\delta$ にも非定数の平方因子が生じる。これで形式的にも被約既約性が従う。

最後に、直接微分によって
\begin{equation}\label{eq:four-saito-identity-v24}
 (\partial_u\delta,\partial_v\delta,\partial_w\delta,\partial_z\delta)(-C)
 =\delta\,(w/2,\,0,\,-8,\,3z/2)
\end{equation}
を得る。$\det(-C)=\delta$ と被約性にSaitoの判定を適用すればよい
\cite[Propositions 1.3, 1.5]{BEvB}。
\end{proof}

自由性そのものは、最小versal有限平坦写像に関する一般定理の適用でもある
\cite[Theorems 3.1, 3.4]{BEvB}。
ここでは本研究のcorank二の被覆について、trace、different、最小表示、
対数的ベクトル場を同じ四次行列で具体化した。
式\eqref{eq:four-quartic-index-v24}の因子 $u^2$ は、$u=0$ で
$X$ 一つでは有限代数の基底を与えられなくなることに対応する。
これを除かない四次判別式をbranch divisorの被約方程式とすることはできない。

\subsection{分岐因子の正規化}

\begin{theorem}\label{thm:four-branch-normalization-v24}
有限射
\[
 V(q)\longrightarrow\Delta
\]
は正規化写像である。従ってbranch divisorの正規化は三次元普通二重点である。
さらに $N=T^1_{\mathscr D/R}$ を有限 $R$-加群と見れば
\begin{equation}\label{eq:four-fast-normalization-v24}
 \End_R(N)\simeq\mathscr D/(q)
 \simeq\mathbb C[[\alpha,\beta,\gamma,\eta]]/(\alpha\beta+\gamma\eta).
\end{equation}
実際の十三変数の基底では、これらの環と写像に九変数の形式円板を加える。
\end{theorem}

\begin{proof}
被約既約な素因子 $(\delta)$ で $R$ を局所化すると離散付値環となる。
有限自由加群 $\mathscr D$ 上の $q$ 倍写像の行列式は $-8\delta$ なので、
その余核 $(\mathscr D/(q))_{(\delta)}$ の長さは一である。
この余核は $\Frac(R/(\delta))$ 上有限な整域、従って体であり、
その長さは $\Frac(\mathscr D/(q))$ の拡大次数に等しい。
従って有限射 $V(q)\to\Delta$ は形式環上でも一般次数一である。
$V(q)$ は正規であるから、この有限双有理射は正規化写像である。

$\det J=q$ と余因子行列の恒等式から $qN=0$ である。
さらに余因子行列は
\[
 N\longrightarrow(\mathscr D/(q))^2,\qquad
 [a]\longmapsto\operatorname{adj}(J)a\bmod q
\]
という単射を与える。実際、$\operatorname{adj}(J)a=qb$ なら
$q a=qJb$ となり、$\mathscr D$ が整域なので $a=Jb$ である。
$J$ は $\Frac(\mathscr D/(q))$ 上で階数一を持つため、
$N$ は $\mathscr D/(q)$ 上の階数一のtorsion-free加群となる。
従って $R/(\delta)$ 上でも階数一かつtorsion-freeであり、
その自己準同型環は共通の関数体の部分環である。
乗法によって $\mathscr D/(q)\subset\End_R(N)$ があり、
他方 $\End_R(N)$ は $R/(\delta)$ 上有限なので、その全ての元は整である。
既に $\mathscr D/(q)$ が整閉包と分かったので、逆包含も従う。
\end{proof}

式\eqref{eq:four-fast-normalization-v24}は、この例における
接コホモロジー加群からの正規化の具体的な実現である。
この加群 $N$ は、次節以降の元の十一点族のpotential障害加群 $M$ とは異なる。
検証器\nolinkurl{four_cover_discriminant_v24.py}は、
乗法関係、trace判別式、Schur消去、式\eqref{eq:four-saito-identity-v24}、
normと四次判別式の係数、foldの横断微分を有理数上で再計算する。
